\chapter{Related Work}
\label{chap:related_work}

This chapter summarises previous work related to this thesis. Specifically, it covers previous achievements and findings of the \ac{LDCI2027} project from a control perspective, as well as general research on \ac{RL}-based \ac{HEV} control strategies.

\section{Previous Light-Duty CI-ICE 2027 Progress}
The \ac{LDCI2027} project focuses on developing advanced operating strategies to minimise \ac{NOx} emissions in light-duty hybrid diesel vehicles, using \ac{AI} and \ac{RL} to support compliance with strict emission regulations such as Euro 6 and Euro 7 \cite{european_union_regulation_2007, european_union_regulation_2024}. Previous progress within this project can be divided into two development phases: the initial parametrisation and testing of control algorithms in MATLAB/Simulink, and the subsequent transition to advanced recurrent \ac{DRL} architectures in Python to address partial observability and computational bottlenecks.

\subsection{Initial Modelling and Base Control Strategies}
\label{subsec:initial_modelling_base_control}
Early work in the project focused on replacing computationally demanding physics-based simulations with data-driven surrogate models \cite{universitat_politecnica_de_catalunya_development_2025}. The authors used \ac{LSTM} \ac{NN}s provided by the \ac{IFS} to predict the vehicle's dynamic behaviour, particularly the battery's \ac{SOC} and tailpipe \ac{NOx} emissions \cite{cuberos_munoz_developing_2025}.

To control these models, the initial methodologies used pre-processing frameworks based on lookup tables and brake-specific fuel consumption curves. This allowed the controller to map high-level actions, such as the power-split factor, to specific engine parameters, including fuel injection mass \cite{cuberos_munoz_developing_2025}.

Using this environment, two control strategies were initially tested \cite{cuberos_munoz_developing_2025}:

\begin{itemize}
  \item \textbf{Nonlinear Model Predictive Control (NLMPC).} An NLMPC was implemented to balance \ac{SOC} tracking and \ac{NOx} minimisation. However, the simulations showed that although the controller could maintain the \ac{SOC} reference, it struggled significantly with emissions, resulting in \ac{NOx} levels nearly five times higher than the target constraints.
  \item \textbf{\ac{SAC} Agent.} An early \ac{RL} approach was tested using a \ac{SAC} algorithm. This implementation highlighted the critical challenge of multi-objective reward scaling. The penalty for \ac{NOx} emissions was found to be orders of magnitude larger than the \ac{SOC} reward, which strongly biased the agent. The \ac{SAC} algorithm ultimately failed to converge within the tested episodes, demonstrating the need for more robust training frameworks and improved reward engineering.
\end{itemize}

\subsection{Model Predictive Control}
Recognising the limitations of the initial implementations, later progress shifted towards a Python-based framework to develop an incremental \ac{NN}-based controller using supervised learning. This phase introduced several methodological improvements \cite{universitat_politecnica_de_catalunya_development_2025}:

\begin{itemize}
  \item \textbf{Incremental action prediction.} Instead of predicting absolute control values, the new controller was redesigned to predict incremental variations ($\Delta$s) using stateful \ac{LSTM}s with a multilayer perceptron head. This approach simplified learning, prevented gradient explosion, and provided more stable control.
  \item \textbf{Differentiable physical constraints.} A major challenge was that the \ac{ICE} model is physically valid only at engine speeds above 900\,rpm; below this threshold, the fuel mass must be strictly zero. To work with this restriction, the authors implemented ``soft activation layers'' in the form of sigmoid gates. This allowed the network to learn differentiably how to shut off fuel injection when the engine speed dropped below the threshold.
\end{itemize}

\subsection{Recurrent Deep Reinforcement Learning}
Despite the improvements achieved with the supervised predictive controller, performance degraded at high speeds and under extreme operating conditions \cite{universitat_politecnica_de_catalunya_development_2025}. Furthermore, a paradigm shift was required: because the internal memory states of the surrogate \ac{LSTM}s cannot be measured in a real vehicle, they must remain hidden from the controller \cite{universitat_politecnica_de_catalunya_development_2025}. These hidden states encode thermal dynamics that influence aftertreatment effectiveness, including the light-off behaviour introduced in \autoref{subsec:thermal_state_complexity}. This strict reliance on measurable physical variables transformed the problem into a \ac{POMDP} \cite{universitat_politecnica_de_catalunya_development_2025}.

To address this issue, the project shifted entirely to a \ac{DRL} framework, specifically selecting a recurrent \ac{TD3} algorithm. This algorithm was chosen because its deterministic policy enables precise setpoint tracking, its off-policy nature improves sample efficiency, and its twin-critic design stabilises continuous control \cite{universitat_politecnica_de_catalunya_development_2025}:

\begin{itemize}
  \item \textbf{Agent architecture.} The architecture used separate \ac{LSTM}s for the actor and the critics to prevent representation bottlenecks and to ensure that the critics evaluated temporal dependencies independently of the policy.
  \item \textbf{Overlapping trajectories.} The researchers implemented an overlapping-window replay buffer with a ``burn-in'' phase to warm up the hidden states before backpropagation. This was necessary because recurrent sampling from the replay buffer must preserve temporal dependencies. During burn-in, a fixed number of initial steps are executed only in forward mode, without propagating gradients back through them.
  \item \textbf{Distributed learning via Ray \cite{liang_rllib_2017}.} Data collection was significantly accelerated using a distributed actor--learner architecture, with environment simulations running in parallel on CPUs while a centralised GPU learner optimised the policy.
\end{itemize}

\subsection{Current Limitations and Open Challenges}
\label{subsec:current_limitations}
Although the established theoretical pipeline and system architecture are advanced and modular, the latest recurrent \ac{TD3} implementation was unable to learn a robust strategy that fully stabilises the vehicle speed at the target value \cite{universitat_politecnica_de_catalunya_development_2025}. The author hypothesises that this is primarily due to a computational sampling bottleneck: simulating the environment with the surrogate \ac{LSTM}s restricts data collection, leaving too few interactions for \ac{DRL} convergence in continuous spaces \cite{universitat_politecnica_de_catalunya_development_2025}. Additionally, potential implementation bugs in the recurrent state propagation or burn-in logic also remain an open challenge \cite{universitat_politecnica_de_catalunya_development_2025}.

The two issues identified in \autoref{subsec:initial_modelling_base_control}, together with the open challenge of the sampling bottleneck, motivate the methodology of this thesis. The emissions shortfall is addressed through the Phase~2 \ac{NOx} penalty and charge-sustaining reward design; the reward-scaling and convergence problem is addressed through reward normalisation, staged training, and reward-weight optimisation; and the sampling bottleneck is addressed by using a newer version of the surrogate \ac{LSTM} models (see \autoref{chap:methodology} and \autoref{sec:impl_lstms}).

\begingroup
\renewcommand{\sectionmark}[1]{\markright{\thesection\ State of the Art -- RL-Based HEV Control}}
\section{State of the Art -- Reinforcement-Learning-Based Hybrid-Electric-Vehicle Control}
\endgroup

The state of the art in \ac{RL}-based \ac{HEV} control is broad, covering different objectives, challenges, approaches, and algorithms \cite{wang_comparative_2023,liu_intelligent_2021, liu_deep_2022}. Many studies address individual aspects of this multi-objective problem, such as avoiding local minima through exploration \cite{zhou_self-learning_2022}, handling hybrid action spaces with separate networks for engine-throttle control and gear shifting \cite{tang_double_2022}, or predicting vehicle velocity when future route knowledge is unavailable \cite{liu_reinforcement_2017}. Because these use cases do not fully align with the present thesis, they mainly provide methodological components that can be adapted here. The most relevant studies are discussed in the following.

\subsection{Continuous Control with Actor-Critic Architectures}
Traditional value-based \ac{RL} methods, such as Q-learning, struggle with the high-dimensional and continuous control spaces inherent in \ac{HEV} energy management. These methods rely on evaluating discrete action bins, which makes them susceptible to discretisation errors and the curse of dimensionality when applied to continuous variables. To overcome these limitations, the state of the art increasingly relies on actor--critic architectures, as described in \autoref{subsec:actor-critic-architectures} \cite{tang_double_2022, zhou_novel_2021}.

Recent \ac{HEV} literature shows a clear transition towards state-of-the-art actor--critic algorithms, primarily Deep Deterministic Policy Gradient (DDPG), \ac{SAC}, and \ac{TD3}. For example, Tresca et al. \cite{tresca_cutting-edge_2025} implemented a \ac{SAC} agent to manage a plug-in \ac{HEV}, highlighting the algorithm's entropy maximisation mechanism, which naturally enhances training stability and exploration. To prevent the neural networks from overfitting to a single driving route, they introduced a multi-cycle training strategy across 72 different mission profiles, showing that training over an ensemble of cycles can substantially improve controller generalisation and robustness under unseen real-world conditions. Similarly, Zhou et al. \cite{zhou_novel_2021} addressed the inherent instability and overestimation bias of standard DDPG agents by employing an improved \ac{TD3} algorithm. Their \ac{TD3}-based \ac{EMS} effectively smoothed the continuous torque commands, leading to faster convergence and more stable global performance than traditional value-based and policy-based methods.

\subsection{Multi-Objective Staged Course Learning}
Formulating an \ac{EMS} for intelligent \acp{HEV} requires balancing multiple conflicting objectives simultaneously, such as minimising fuel consumption, reducing tailpipe emissions, maintaining battery \ac{SOC}, and ensuring passenger comfort \cite{wang_energy_2025}. Aggregating these competing goals into a single static reward function can destabilise the \ac{RL} process. The agent may receive interfering reward signals, leading to premature convergence or trapping the policy in local optima, thereby preventing optimisation of the intended objective.

To address the complexities of high-dimensional multi-objective control, Wang and Wang \cite{wang_energy_2025} proposed a multi-stage ``course learning'' mechanism integrated within a \ac{TD3} framework for managing both \ac{ACC} and \ac{EMS}. Drawing inspiration from progressive human learning, this approach breaks the complex control problem down into sequential sub-stages with different focuses. In the initial training stage, the agent is strictly guided to satisfy primary safety prerequisites, receiving rewards exclusively for maintaining a safe car-following distance and basic speed tracking. Once the agent successfully learns to keep the speed-tracking error within an acceptable boundary, the training automatically transitions to a secondary stage that refines the control-action quality. In this subsequent phase, variable virtual goals are introduced and energy-saving rewards are amplified, allowing the agent to optimise secondary objectives, such as fuel economy and power-source durability, while retaining the foundational safety rules.

\subsection{Domain Knowledge Integration}
While model-free \ac{DRL} algorithms offer powerful self-learning capabilities, their initial trial-and-error exploration phases can generate irrational, physically infeasible, or otherwise unfavourable control actions. Relying purely on the neural network to autonomously discover and represent the physical constraints of the \ac{ICE}, electric motors, and battery pack is sample-inefficient and may not sufficiently cover the real control problem, especially when component-health-preserving measures are not correctly encoded in the \ac{RL} loop. Consequently, recent research trends focus on actively embedding domain knowledge and heuristic rules directly into the \ac{DRL} loop to constrain the action space.

Zhou et al. \cite{zhou_novel_2021} designed a heuristic rule-based local controller that acts as a dynamic action-space regulator within their \ac{TD3}-based \ac{EMS}. By computing the instantaneous allowable engine and motor torques before neural-network exploration, the local controller filters out infeasible or irrelevant torque allocations, reducing the search space and providing a ``warm start'' that accelerates convergence. Furthermore, Tang et al. \cite{tang_double_2022} showed that \ac{DRL} algorithms generally struggle to learn intermittent and discontinuous control commands, such as gear shifting or mode switching, because \ac{DRL} networks output continuous actions at every time step. To address this limitation, they embedded a rule-based engine start--stop strategy into their Double \ac{DRL} framework. By assigning the intermittent decision of turning the engine on or off to a deterministic rule-based system, the continuous \ac{RL} agent could focus on optimising the throttle, resulting in fuel-economy improvements over traditional \ac{DP} benchmarks.
